\documentclass[11pt]{article}
\usepackage[margin=1in]{geometry}
\usepackage{amsmath,amssymb,booktabs,array,graphicx,hyperref,microtype,xcolor}
\usepackage[T1]{fontenc}
\usepackage[section]{placeins}

\graphicspath{{../results/06_window_parameter_tuning/}}
\hypersetup{
  colorlinks=true,
  linkcolor=blue!55!black,
  urlcolor=blue!65!black,
  citecolor=blue!55!black
}
\setlength{\parindent}{0pt}
\setlength{\parskip}{0.62em}
\newcommand{\Rfixed}{R^2_{\mathrm{fixed}}}
\newcommand{\Rlocal}{R^2_{\mathrm{local}}}

\title{Systematic Time-Window, Delay, and Rank Verification\\
for Local DMD of \texttt{mouse02\_sleep}}
\author{Reproducible development-only analysis report}
\date{17 July 2026}

\begin{document}
\maketitle

\begin{abstract}
This report determines what time window and calibration coefficient were used
in the preceding analysis, then tests a standard local dynamic mode
decomposition (DMD) tuning pipeline on one NREM recording from
\texttt{mouse02\_sleep}.  The preceding code selected a 196.08-s raw block but
fitted DMD to only 117.65 s.  The corrected analysis uses PyDMD 2025.8.1,
four-frame event-rate bins, a fixed acquisition-1 observation basis, common
chronological forecast origins, and a joint scan over 20--180-s fitting
windows, five delay depths, and seven explicit ranks.  Two prediction
coefficients are kept deliberately separate.  The primary
\(\Rfixed\) uses the same acquisition-1 mean reference for every window and
therefore permits a fair comparison across \(W\); the auxiliary \(\Rlocal\)
measures only the increment beyond each window's own past mean.

No guarded configuration achieved positive \(\Rfixed\), and none passed the
technical development gate.  The lowest comparable loss occurred at
\(W=180\) s, \(d=1\), \(r=12\), with
\(\Rfixed=-0.004833\) and \(\Rlocal=0.004651\).  Because there was no
gate-passing pool, a one-standard-error parsimony calculation produced only a
\emph{fallback}: \(W=180\) s, \(d=1\), \(r=2\), with
\(\Rfixed=-0.005171\) and \(\Rlocal=0.004224\).  The 45-s, rank-4 model has
the largest local increment, \(\Rlocal=0.008932\), but its comparable score is
\(\Rfixed=-0.042152\).  Thus the fair forecast screen favors longer windows
only in the limited sense that they lose less; it does not verify a 180-s
model.

A damping-aware spectral audit found no eligible oscillatory pair in either
the 45-s local-increment candidate or the 180-s fallback.  Longer delayed
models sometimes produced eligible signatures, but the available origins
were too overlapping to verify recurrence.  A separately refitted,
development-internal Residual-DMD-style diagnostic gave a globally
synchronized surrogate \(p=0.13\), unadjusted for candidate selection.
Training reconstructions explained essentially none of the all-neuron
variance.  All optimized-DMD fits emitted convergence warnings.  The outer
41.83-s tail consequently remains score-locked and unscored.  Neither 45 s
nor 180 s is a verified window for brain-state mode tracking.
\end{abstract}

\section{Direct answers and interpretation}

\subsection{What was the previous DMD window?}

Three durations had been conflated.  The selected raw block contained 1500
native frames, or
\[
\frac{1500}{7.65}=196.08\ \mathrm{s}.
\]
Four-frame binning produced 375 samples.  The PyDMD fit used only the first
225 bins,
\[
225\frac{4}{7.65}=117.65\ \mathrm{s}.
\]
The next and final 75-bin blocks were each 39.22 s.  Therefore 117.65 s,
not 196.08 s, was the preceding DMD fitting duration.

\subsection{What was calibration \texorpdfstring{\(R^2\)}{R2}?}

The earlier quantity was a held-out one-step prediction coefficient pooled
over neuron--time entries,
\[
R^2_{\mathrm{cal}}
=1-\frac{\sum_{i,t}(y_{i,t}-\widehat y_{i,t})^2}
{\sum_{i,t}(y_{i,t}-\overline x_{i,\mathrm{train}})^2}.
\]
It was not probabilistic calibration, autonomous trajectory reconstruction,
a Residual-DMD statistic, or evidence that a dynamic mode was stable.  The
old delay comparison also allowed the automatic SVD rank to change with
delay, so it did not isolate delay depth.  This report avoids the ambiguous
term \emph{calibration \(R^2\)} and instead defines a comparable
fixed-reference prediction coefficient and a separate local-reference
increment.

\subsection{Should the fitting window simply be shorter?}

No general rule makes the shortest window preferable.  A short window may
reduce mixing between changing dynamical regimes, but it also gives fewer
regression transitions, a noisier local mean, weaker rank support, and poorer
resolution of slow oscillations.  A long window reverses those advantages and
disadvantages.  For \(N_W\) binned samples and delay \(d\), the number of
snapshot transitions is \(N_W-d\), the explicit delay memory is
\((d-1)\Delta t\), and the delayed ambient dimension is \(qd\).

The present data do not support the claim that a short window is better.  The
fair cross-window loss decreases toward 180 s, whereas the small improvement
relative to a window-specific local mean peaks at 45 s.  Every comparable
score nevertheless remains negative.  Online and windowed DMD likewise treat
locality as a parameter to validate rather than a universal optimum
\cite{windoweddmd,lapo}.  The scientifically honest conclusion is therefore
not that 180 s works, but that no tested window currently works.

\section{Data, observable, and chronological design}

\subsection{One dataset and one state}

Only \texttt{mouse02\_sleep.mat} and raw state code 1 (NREM) were used.  No
paper-specific \texttt{used\_frame} mask entered this analysis.  Every bin,
delay vector, fitting window, and target remained within one NREM run and one
acquisition.  This state restriction prevents an explicit scored state
transition inside a window, but it does not prove that a single stationary
dynamical system generated the complete NREM interval.

Four native deconvolved-event samples were summed and divided by
\[
\Delta t=\frac{4}{7.65}=0.522876\ \mathrm{s}.
\]
The resulting observable is event mass per second.  It is a continuous-valued
rate proxy constructed from event data, not a calibrated spike rate.

\subsection{Fixed neuron transformation and observation basis}

Acquisition-1 NREM bouts supplied the fixed neuron transformation.  Of 6,574
neurons, 6,570 had finite nonzero acquisition-1 RMS.  For eligible neuron
\(i\),
\[
x_i(t)=
\frac{r_i(t)-\mu_{i,\mathrm{acq1}}}
{\mathrm{RMS}_{i,\mathrm{acq1}}+10^{-3}}.
\]
A POD basis \(B_q\) was learned only from these standardized acquisition-1
NREM samples.  All development targets came from acquisition 2.  For a
candidate fitting window, the DMD coordinates were
\[
z(t)=B_q^\top\{x(t)-\overline x_W\},
\qquad
\widehat x(t+h)=B_q\widehat z(t+h)+\overline x_W,
\]
where \(\overline x_W\) is computed from the past fitting window only.
Forecast loss was evaluated in the complete 6,570-neuron norm.  An exact
orthonormal-projection identity retained error outside
\(\operatorname{col}(B_q)\), so the sweep did not silently score only the
retained POD coordinates.

\subsection{Development and outer geometry}

The longest acquisition-2 NREM bout contained 516 complete rate bins, or
269.80 s.  Eight common development origins were separated by 12 bins; their
eight-bin future target blocks did not overlap, although their fitting
contexts did.  The final 80 bins, or 41.83 s, were assigned to five possible
outer origins.  The corrected run defined this outer interval by time and
state label but did not display or score its neural activity.

An earlier exploratory draft had displayed that activity.  The current report
therefore calls the interval \emph{score-locked and unscored}, not pristine or
untouched.  It is not opened unless the complete development gate passes.

\begin{figure}[p]
\centering
\includegraphics[width=0.98\linewidth]{00_chronological_validation_geometry.png}
\caption{Acquisition-1 representation bouts and the acquisition-2
development/outer geometry.  The orange outer interval is defined, but its
activity and scores remain closed.}
\label{fig:geometry}
\end{figure}

\section{Two prediction references that answer different questions}

Let \(o\) index a forecast origin and
\(\mathcal H=\{1,2,4,8\}\) the scored horizons in rate bins.  The primary loss
is the equal-weight mean
\[
\mathcal L_{\mathrm{fixed}}
=
\operatorname{mean}_{o,\,h\in\mathcal H}
\frac{\mathrm{SSE}_{\mathrm{DMD}}(o,h)}
{\mathrm{SSE}_{\mu_{\mathrm{acq1}}}(o,h)},
\qquad
\Rfixed=1-\mathcal L_{\mathrm{fixed}}.
\]
After acquisition-1 standardization,
\(\mu_{\mathrm{acq1}}=0\).  At a given origin and horizon the denominator is
identical for every \(W,d,r\), which makes \(\Rfixed\) the appropriate
cross-window selection metric.

The auxiliary local increment is
\[
\mathcal L_{\mathrm{local}}
=
\operatorname{mean}_{o,\,h\in\mathcal H}
\frac{\mathrm{SSE}_{\mathrm{DMD}}(o,h)}
{\mathrm{SSE}_{\overline x_W}(o,h)},
\qquad
\Rlocal=1-\mathcal L_{\mathrm{local}}.
\]
It asks whether the autonomous DMD forecast improves on the same window's
past-only constant prediction.  This is scientifically useful, but its
denominator changes with \(W\), so it cannot fairly select a window.

\begin{table}[htbp]
\centering
\caption{Mean denominator over the common development origins and horizons.
The fixed denominator is constant across \(W\); the local denominator is not.}
\label{tab:denominators}
\begin{tabular}{rrr}
\toprule
\(W\) (s) & fixed-reference SSE & local-mean SSE \\
\midrule
20  & 10880.72 & 11375.86 \\
30  & 10880.72 & 11238.00 \\
45  & 10880.72 & 11112.39 \\
60  & 10880.72 & 11026.05 \\
90  & 10880.72 & 10951.64 \\
120 & 10880.72 & 10914.67 \\
180 & 10880.72 & 10892.08 \\
\bottomrule
\end{tabular}
\end{table}

The short-window local mean is, on average, worse than the fixed
acquisition-1 mean.  A model can consequently have a small positive
\(\Rlocal\) while remaining negative under \(\Rfixed\).  That is exactly what
occurs here.  Reporting only \(\Rlocal\) would make the 45-s peak look like a
cross-window optimum when it is not.

\section{Representation-capacity audit}

The original \(q=64\) POD basis retained only 22.03\% of standardized
acquisition-1 energy.  Before fitting DMD, the corrected analysis computed an
instantaneous projection oracle on the development targets.  For each future
truth, the oracle used the best possible prediction in
\(\overline x_W+\operatorname{col}(B_q)\).  It is future-informed and is
therefore not forecast performance.  It asks only whether the observation
space makes the declared effect-size target attainable.

\begin{table}[htbp]
\centering
\caption{Nested POD capacity.  The two oracle columns are maxima over the
tested windows and use the indicated prediction reference.}
\label{tab:oracle}
\begin{tabular}{rrrr}
\toprule
\(q\) & retained energy & max fixed oracle \(R^2\) & max local oracle \(R^2\) \\
\midrule
64  & 22.03\% & 0.02022 & 0.03620 \\
128 & 34.76\% & 0.03155 & 0.04739 \\
256 & 52.74\% & 0.05260 & 0.06811 \\
\bottomrule
\end{tabular}
\end{table}

The 64- and 128-dimensional representations could not reach the declared
\(R^2=0.05\) target even with future information.  The corrected run fixed
\(q=256\), the smallest tested dimension that made the target attainable.
The practical effect target was prespecified on \(\Rlocal\); the fixed oracle
is additionally reported because \(\Rfixed\) controls cross-window selection.
Every fitted score was checked programmatically to remain below the
corresponding oracle ceiling.

\begin{figure}[htbp]
\centering
\includegraphics[width=0.84\linewidth]{01b_pod_projection_capacity.png}
\caption{Future-informed development projection ceilings.  These curves
measure observation-space capacity, not DMD prediction.}
\label{fig:oracle}
\end{figure}

\FloatBarrier
\section{Locked window, delay, and rank tuning pipeline}

\subsection{Candidate grid and support guard}

The candidate grid was
\[
W\in\{20,30,45,60,90,120,180\}\ \mathrm{s},
\qquad
d\in\{1,2,4,8,16\},
\]
\[
r\in\{2,4,8,12,16,24,32\}.
\]
Rounded fitting durations were 19.869, 29.804, 44.967, 60.131, 89.935,
120.261, and 179.869 s.  Delay memories were 0, 0.523, 1.569, 3.660, and
7.843 s.  A candidate entered selection only if
\[
\frac{N_W-d}{r}\geq 8.
\]
This eight-transitions-per-rank requirement is a transparent project
heuristic, not a DMD theorem.  It retained 143 of 245 displayed combinations.

\begin{figure}[htbp]
\centering
\includegraphics[width=0.98\linewidth]{03_grid_support_and_frequency_tradeoff.png}
\caption{Regression support, delay memory, and rough frequency-resolution
trade-offs over the locked grid.}
\label{fig:grid}
\end{figure}

\subsection{Primary estimator and autonomous scoring}

The primary estimator was PyDMD's \texttt{HankelDMD} with explicit
\texttt{svd\_rank}, \texttt{exact=True}, and the candidate delay \(d\).
TLSQ rank was zero; optimized amplitudes, forward--backward correction,
rescaling, and Tikhonov regularization were disabled.  PyDMD exposes these as
separate choices \cite{hankeldmd}; the primary screen changed only \(W,d,r\)
so their effects remained interpretable.

At each origin the complete past window was fitted once.  The last observed
delay state initialized one autonomous eight-bin forecast.  No future
observation was inserted after initialization.  Persistence and a stable
neuronwise AR(1) were additional dynamic baselines.  The implementation also
contains a noise-free autonomous-forecast unit check and an exact all-neuron
POD-score identity check.

\subsection{Technical and practical gates}

A technically successful configuration required:
\begin{enumerate}
\item positive mean \(\Rfixed\) and \(\Rlocal\);
\item positive fixed and local scores separately at \(h=1\) and \(h=8\);
\item positive-origin fractions of at least 0.75 for both references overall
and at the endpoint horizons;
\item positive mean skill against persistence and AR(1), with at least 0.75
of origins beating both;
\item finite numerical diagnostics and no unexpected warning.
\end{enumerate}
The practical gate additionally required \(\Rlocal\geq0.05\), stable
selection, recurrent modes, and the separate residual diagnostic.  Repeated
rolling origins and multiple horizons follow general out-of-sample forecast
evaluation practice \cite{tashman}; the numerical thresholds remain project
heuristics.

A delete-two-origin jackknife supplied a descriptive one-standard-error
parsimony calculation.  The fitting contexts overlap, so this is not a formal
independent-sample standard error.  Most importantly, when no configuration
passes the technical gate, its output is labelled a \emph{fallback}, not a
selected or verified model.

\section{Forecast results: longer loses less, but no window succeeds}

All 1,144 guarded exact-Hankel development fits returned and none failed.
There were zero technically passing configurations.  All 143 configurations
had negative \(\Rfixed\), while 98 had positive \(\Rlocal\).  The latter
counts incremental improvement over a window-dependent and often inferior
local-mean baseline; it is not evidence that 98 models forecast usefully.

\begin{table}[htbp]
\centering
\caption{Lowest comparable loss within each \(W\).  The table reports both
references because they answer different questions.}
\label{tab:window}
\begin{tabular}{rrrrrrr}
\toprule
\(W\) (s) & \(d\) & \(r\) & \(\Rfixed\) & \(\Rlocal\) &
skill--persist. & skill--AR(1) \\
\midrule
20  & 1 & 2  & -0.084123 & 0.006792 & 0.4871 & 0.2026 \\
30  & 1 & 4  & -0.060204 & 0.007092 & 0.4965 & 0.1839 \\
45  & 1 & 4  & -0.042152 & 0.008932 & 0.5036 & 0.1604 \\
60  & 1 & 12 & -0.028138 & 0.007543 & 0.5087 & 0.1273 \\
90  & 1 & 2  & -0.015195 & 0.005500 & 0.5135 & 0.1114 \\
120 & 4 & 2  & -0.009526 & 0.004363 & 0.5155 & 0.0867 \\
180 & 1 & 12 & -0.004833 & 0.004651 & 0.5179 & 0.0287 \\
\bottomrule
\end{tabular}
\end{table}

The raw comparable minimum is \(W=180\) s, \(d=1\), \(r=12\).  Its positive
skill against persistence and AR(1) does not contradict its negative
\(\Rfixed\): both dynamic baselines are worse than the fixed acquisition mean.
The rank-12 model is therefore the \emph{least poor comparable forecast}, not
a successful model.

No gate-passing pool existed.  Applying the descriptive one-SE rule to the
complete guarded grid produced \(W=180\) s, \(d=1\), \(r=2\).  This result is
recorded as \path{fallback_no_gate_passing_configuration}.  Its
horizon-specific scores are:

\begin{table}[htbp]
\centering
\caption{Development scores for the 180-s, \(d=1,r=2\) fallback.}
\label{tab:fallback-horizon}
\begin{tabular}{rrrr}
\toprule
horizon (bins) & horizon (s) & \(\Rfixed\) & \(\Rlocal\) \\
\midrule
1 & 0.523 &  0.000816 &  0.008923 \\
2 & 1.046 & -0.008724 &  0.002350 \\
4 & 2.092 & -0.010630 &  0.005676 \\
8 & 4.183 & -0.002145 & -0.000051 \\
\bottomrule
\end{tabular}
\end{table}

The fallback fails both the mean common-reference requirement and the
eight-bin local-reference endpoint.  It must not be called the final verified
model.

\begin{figure}[p]
\centering
\includegraphics[width=0.98\linewidth]{04_development_window_delay_heatmaps.png}
\caption{Fixed-reference and local-reference development scores.  The
fixed-reference panels permit comparison across \(W\); the local panels show
only the increment beyond each window's own past mean.}
\label{fig:heatmaps}
\end{figure}

\begin{figure}[htbp]
\centering
\includegraphics[width=0.98\linewidth]{05_selected_development_scores.png}
\caption{Origin and horizon diagnostics for the one-SE fallback.  A displayed
fallback is not a gate-passing selection.}
\label{fig:selected}
\end{figure}

\subsection{Delay does not rescue short windows}

\begin{table}[htbp]
\centering
\caption{Best comparable result for each delay over all guarded \(W,r\).}
\label{tab:delay}
\begin{tabular}{rrrrr}
\toprule
\(d\) & \(W\) (s) & \(r\) & \(\Rfixed\) & \(\Rlocal\) \\
\midrule
1  & 180 & 12 & -0.004833 & 0.004651 \\
2  & 180 & 2  & -0.006295 & 0.003090 \\
4  & 180 & 12 & -0.006835 & 0.002602 \\
8  & 180 & 2  & -0.007278 & 0.002230 \\
16 & 180 & 2  & -0.008596 & 0.000976 \\
\bottomrule
\end{tabular}
\end{table}

The hypothesized pattern that a short \(W\) needs a large \(d\) is not
supported for this observable.  Delay embedding may still be necessary for
other observables and systems; PyDMD's documented workflow recommends
examining singular support and validation performance rather than assuming a
delay depth \cite{pydmdtutorial}.

\subsection{Selection instability}

Deleting consecutive origin pairs produced fallback one-SE tuples
\[
(W,d,r)\in
\{(120,4,2),(180,1,2),(180,1,2),(180,1,2)\}.
\]
The full fallback recurred in three of four screens, but none of the four
reduced screens had any gate-passing pool.  The 0.75 tuple match rate is
therefore not successful scientific selection stability.

\section{Numerical and noise-aware checks}

\subsection{Exact-Hankel warnings and conditioning}

Of 1,144 exact fits, 864 emitted PyDMD's input-condition warning.  All were
classified as expected rank-truncation warnings and none was unexpected.
The warnings occurred for the shorter matrices; the 180-s finalists emitted
none.  No exact fit failed.

For the 180-s, \(d=1,r=12\) comparable minimum, the median retained condition
number \(\sigma_1/\sigma_{12}\) was 4.171, while the median boundary ratio
\(\sigma_{12}/\sigma_{13}\) was only 1.016.  For the rank-2 fallback, the
corresponding medians were 1.909 and 1.300.  Full operator-pair and centered
input condition numbers were approximately 99.  The explicit truncations are
numerically manageable, but the weak rank-12 boundary does not support
interpreting 12 as a sharply identified intrinsic dimension.

\subsection{Optimized/BOP-DMD sensitivity}

PyDMD's noisy-data tutorial demonstrates delay coordinates with optimized and
bagged optimized DMD \cite{pydmdtutorial,bopdmd}.  A locked shortlist therefore
tested optimized DMD with conjugate-pair constraints, with and without an
additional stability constraint.  In PyDMD, \texttt{num\_trials=0} invokes
optimized DMD without bootstrap bagging.

There were 16 estimator/configuration variants and 128 origin-level fits.
Every origin fit emitted a convergence warning.  Their returned objective
values are not treated as valid performance estimates, and no
convergence-clean optimized candidate was available to promote.  Variable
projection can be sensitive to nonlinear initialization and local solutions
\cite{askham}; a returned object alone is not evidence of successful
optimization.

\begin{figure}[htbp]
\centering
\includegraphics[width=0.90\linewidth]{06_bopdmd_shortlist_comparison.png}
\caption{Exact-Hankel shortlist performance.  Nonconverged optimized-DMD
scores are deliberately omitted from interpretation.}
\label{fig:bop}
\end{figure}

\subsection{Raw-rate inner-product sensitivity}

Neuronwise RMS scaling gives low- and high-rate neurons similar weight.  The
complete grid was therefore repeated in raw-centered event-rate units with a
separate fixed \(q=256\) basis, which retained 58.36\% of acquisition-1
raw-rate energy.

The fair raw-rate minimum was \(W=180\) s, \(d=1\), \(r=4\), with
\(\Rfixed=-0.003086\) and \(\Rlocal=0.005302\).  The largest raw local
increment was again at \(W=45\) s, \(d=1\), \(r=4\):
\(\Rlocal=0.010190\), but \(\Rfixed=-0.029063\).  No raw-rate configuration
passed.  The sensitivity analysis reproduces the distinction between a
long-window comparable minimum and a small 45-s local increment; it does not
validate either model.

\begin{figure}[htbp]
\centering
\includegraphics[width=0.96\linewidth]{09_raw_vs_standardized_window_sensitivity.png}
\caption{Raw-rate and standardized scoring.  Comparable fixed-reference
selection and window-specific local increments are shown separately.}
\label{fig:raw}
\end{figure}

\FloatBarrier
\section{Within-window locality is still an approximation}

State-label homogeneity prevents an explicit NREM-to-other-state transition
inside a scored window; it does not establish stationarity.  A descriptive
diagnostic split the first eight fixed POD coordinates into chronological
halves and compared their means, covariances, ridge-stabilized one-step
operators, and cross-half one-step transfer.

\begin{table}[htbp]
\centering
\caption{Descriptive split-half locality diagnostics.}
\label{tab:drift}
\begin{tabular}{rrrrrr}
\toprule
\(W\) (s) & mean shift & covariance drift & operator drift &
transfer \(R^2\) & positive fraction \\
\midrule
20  & 0.646 & 1.439 & 1.378 & -1.8885 & 0.000 \\
30  & 0.558 & 1.189 & 1.306 & -0.2696 & 0.625 \\
45  & 0.471 & 0.983 & 1.207 &  0.1805 & 0.750 \\
60  & 0.429 & 0.933 & 1.081 &  0.3917 & 1.000 \\
90  & 0.300 & 0.601 & 1.089 &  0.5133 & 1.000 \\
120 & 0.343 & 0.577 & 0.941 &  0.5501 & 1.000 \\
180 & 0.187 & 0.318 & 0.776 &  0.5627 & 1.000 \\
\bottomrule
\end{tabular}
\end{table}

These values do not show that short windows are more stationary.  They also
do not certify the 180-s window: a longer half contains more fitting samples,
so improved transfer can reflect reduced estimator variance rather than a
more constant underlying system.  The diagnostic is not a gate in this
pilot.  A confirmatory study should prespecify state-preserving block
surrogates or change-point exclusions before using drift as a decision rule.

\begin{figure}[htbp]
\centering
\includegraphics[width=0.98\linewidth]{09b_within_window_drift.png}
\caption{Split-half locality diagnostics in a fixed eight-dimensional space.
They expose, but do not resolve, the stationarity approximation.}
\label{fig:drift}
\end{figure}

\section{Damping-aware mode recurrence}

\subsection{Why frequency alone is insufficient}

Forecast success and spectral adequacy are separate requirements.  A
frequency-resolved eigenvalue can decay so quickly that it does not describe
a recurrent oscillatory component.  For a positive-frequency conjugate pair,
the audit therefore used the actual evidence span
\[
T_{\mathrm{evidence}}=(N_W-d)\Delta t
\]
and required
\[
fT_{\mathrm{evidence}}\geq3
\]
cycles.  It also required at least 1\% of a descriptive fitted-component
Frobenius-power proxy.  This proxy uses per-snapshot least-squares
coefficients; it is not physical mode energy.

For continuous rate
\[
\omega=\frac{\log\lambda}{\Delta t},
\]
the retained amplitude after one oscillatory cycle is
\[
m_{\mathrm{cycle}}
=\exp\left(\frac{\operatorname{Re}\omega}{f}\right).
\]
The pilot persistence threshold was \(m_{\mathrm{cycle}}\geq0.10\): at least
10\% of the amplitude must survive one period.  This transparent threshold is
a project heuristic, not a DMD theorem.

Eligible pairs were matched across origins when their frequencies differed
by no more than one Rayleigh bin, the affinity between their real spatial
mode planes was at least 0.70, and their log per-cycle multipliers differed by
at most one.  A recurrent track had to cover at least six of eight origins
and directly match at least 75\% of at least two available origin pairs whose
fitting contexts shared no more than 25\% of their samples.

\subsection{Observed mode inventory}

Across the deterministic exact-Hankel shortlist there were 179
positive-frequency candidate pairs.  Of these, 160 completed at least three
cycles, 173 passed the 1\% proxy threshold, 46 passed the damping threshold,
and 45 passed all eligibility conditions.  The median per-cycle multiplier
over all candidate pairs was only 0.01006.

The candidate-specific results are decisive:
\begin{itemize}
\item The 45-s, \(d=1,r=4\) local-increment candidate produced nine
positive-frequency pairs across eight origins.  Eight were frequency-resolved
and all nine were salient, but none retained 10\% amplitude for one cycle.
The median multiplier was 0.02929 and the maximum was 0.03984.
\item The 180-s, \(d=1,r=12\) comparable minimum produced 38
positive-frequency pairs.  None passed the damping threshold; the median
multiplier was \(1.69\times10^{-5}\) and the maximum was 0.03687.
\item The 180-s, \(d=1,r=2\) fallback produced no positive-frequency
conjugate pair.
\end{itemize}

Some delayed 60- and 180-s models did produce damping-qualified signatures.
However, every \(W\geq60\) candidate had fewer than two sufficiently
separated origin pairs in this bout.  Adjacent fitting contexts overlapped by
89.57\% at 60 s and 96.51\% at 180 s.  The 45-s geometry offered three
low-overlap origin pairs, but that candidate had no damping-qualified mode.
Consequently no exact-Hankel mode track was verified.  Long-window recurrence
is \emph{geometrically inconclusive}; it is not evidence of recurrence.

\begin{figure}[p]
\centering
\includegraphics[width=0.98\linewidth]{07_modal_stability_and_resdmd_null.png}
\caption{Damping-qualified frequency signatures, spatial-plane matching
versus fitting-context overlap, and the separate internal Residual-DMD-style
diagnostic.}
\label{fig:modes}
\end{figure}

\begin{figure}[htbp]
\centering
\includegraphics[width=0.98\linewidth]{08_audit_candidate_spatial_modes.png}
\caption{Phase-aligned all-neuron maps for the audit candidate at the last
development origin.  They are descriptive fitted maps; no oscillatory mode
from the fallback passed the spectral gate.}
\label{fig:spatial}
\end{figure}

\section{Residual-DMD-style diagnostic with a synchronized null}

PyDMD has no Residual-DMD fitter; its \texttt{RDMD} class denotes Randomized
DMD.  The present calculation therefore refitted a separate rank-\(r\) EDMD
dictionary to the first part of each context and evaluated the published
\(G,A,L\) residual quadratic form on an internal chronological tail
\cite{resdmd}.  It is not a residual attached to the displayed full-window
PyDMD modes.

Because no spectrally qualified finalist existed, the diagnostic audited the
development fallback \(W=180\) s, \(d=1,r=2\).  Each 344-bin context used 258
bins (134.90 s) for dictionary fitting and 86 bins (44.97 s) for quadrature.
The centering mean and dictionary were learned from the fit prefix only.
These tails had nevertheless already participated in model development, so
they are not post-selection confirmation.

The null used 199 fixed-POD-coordinate circular-shift surrogates.  Each
surrogate was generated once over the complete accessible development prefix:
each POD coordinate received one independent circular shift, and the same
shifted realization was then sliced for all eight origins.  This construction
preserves the overlap-induced dependence between origin statistics under the
null.  It corrects the invalid alternative of shifting each overlapping
origin independently.

All eight origin-level plus-one lower-tail \(p\)-values were between 0.135 and
0.830; none was at most 0.05.  The predeclared global statistic was the median
quadrature residual across origins:
\[
T_{\mathrm{obs}}=0.662734.
\]
Twenty-five of 199 synchronized null medians were at or below the observed
value, giving
\[
p=\frac{25+1}{199+1}=0.13.
\]
The diagnostic is exploratory and is not adjusted for choosing \(W,d,r\) from
the same development outcomes.  In a spontaneous stochastic trajectory, the
statistic also contains dynamical variance as well as projection error
\cite{stochasticresdmd}.  It therefore does not verify a mode family.

\section{Can the fitted modes reconstruct the transformed activity?}

Prediction remained the selection criterion, but a claimed mode family
should at least describe the transformed activity used to fit it.  Three
development candidates were therefore audited in-sample using:
\begin{enumerate}
\item reconstruction \(R^2\) in the complete all-neuron norm, including
variance discarded by the POD basis; and
\item overlap between normalized aggregate POD power spectra.
\end{enumerate}
Both diagnostics are optimistic because the same fitting window determines
the modes and their amplitudes.

\begin{table}[htbp]
\centering
\caption{Descriptive training reconstruction and spectral adequacy.}
\label{tab:reconstruction}
\footnotesize
\begin{tabular}{lrrrr}
\toprule
role & \(W,d,r\) & median reconstruction \(R^2\) &
median spectral overlap & peak matches \\
\midrule
comparable minimum & \(180,1,12\) & 0.000065 & 0.798843 & 0/8 \\
one-SE fallback     & \(180,1,2\)  & 0.000029 & 0.641193 & 0/8 \\
largest local increment & \(45,1,4\) & 0.000277 & 0.740397 & 3/8 \\
\bottomrule
\end{tabular}
\end{table}

The normalized spectral overlap is insensitive to total amplitude and can be
moderately high even when the reconstruction explains virtually no
all-neuron variance.  The 180-s observed aggregate peak was not recovered at
any origin.  These fitted low-rank modes do not precisely capture the
transformed neural activity, even in-sample.  The contrast with the roughly
0.06 projection-oracle ceiling shows that this failure is not solely a
consequence of choosing \(q=256\); the fitted low-rank dynamics use almost
none of the available representational capacity.

\begin{figure}[htbp]
\centering
\includegraphics[width=0.96\linewidth]{08b_training_reconstruction_spectral_audit.png}
\caption{Optimistic in-sample all-neuron reconstruction and normalized
spectral overlap.  Spectral-shape overlap must not be mistaken for variance
capture.}
\label{fig:reconstruction}
\end{figure}

\FloatBarrier
\section{Development decision}

\begin{table}[htbp]
\centering
\caption{Frozen development decision.}
\label{tab:decision}
\begin{tabular}{p{0.58\linewidth}cc}
\toprule
criterion & observed & pass \\
\midrule
Positive fixed-reference configuration & 0/143 & no \\
Technical multi-reference forecast gate & 0/143 & no \\
Practical local \(R^2\geq0.05\) & best 0.008932 & no \\
Gate-passing delete-pair stability & no passing pool & no \\
Damping-qualified pair in the fallback & none & no \\
Low-overlap recurrent exact-Hankel track & none verified & no \\
Global synchronized residual diagnostic & \(p=0.13\) & no \\
Raw-rate practical sensitivity & best 0.010190 & no \\
Convergence-clean optimized-DMD fit & 0/128 & no \\
\bottomrule
\end{tabular}
\end{table}

The development result is \textbf{not satisfactory for scientific mode
tracking}.  The negative outcome is not an implementation failure: the
pipeline completed, numerical identities passed, candidate failures and
warnings were disclosed, and the observation-space target was attainable.
Rather, the present observable and estimator did not produce useful
common-reference prediction, stable selection, reconstructive capture, or
verified recurrent modes.

The two window statements must remain distinct:
\begin{itemize}
\item \(W=180\) s minimizes the comparable forecast loss, but that loss is
still worse than the fixed acquisition mean.  The rank-2 one-SE result is a
fallback only.
\item \(W=45\) s maximizes a small increment beyond its own inferior local
mean, but it has strongly negative common-reference performance and no
damping-qualified oscillatory pair.
\end{itemize}
Neither window is verified.  Opening the outer tail to choose another window,
delay, rank, smoother, or mode would spend confirmation data on continued
development, so its score files intentionally contain headers only.

\section{Recommended standard pipeline and next experiment}

The verification yields the following standard procedure:
\begin{enumerate}
\item define state- and acquisition-safe observables in physical units;
\item learn scaling and any fixed spatial basis on separate data;
\item declare one prediction reference that is constant across all windows,
and report window-specific local increments separately;
\item audit whether the observation space can attain the declared effect size
under both references;
\item tune \(W,d,r\) jointly with an explicit transitions-per-rank guard and a
common-rank slice that isolates \(W,d\);
\item use common chronological origins, endpoint-initialized autonomous
multi-horizon forecasts, and fixed-mean, local-mean, persistence, and AR
baselines;
\item disclose conditioning, rank-boundary support, every warning, and every
fit failure rather than treating a returned object as a successful fit;
\item require optimized estimators to converge before comparing their scores;
\item audit in-sample all-neuron reconstruction and spectral shape separately;
\item require frequency resolution, fitted-component salience, and an
explicit damping/lifetime criterion before calling an eigenpair a recurrent
oscillation;
\item establish recurrence across low-overlap windows, preferably separated
bouts or animals, rather than adjacent highly overlapping contexts;
\item generate one synchronized surrogate realization across all overlapping
origins and either fix the candidate in advance or repeat selection inside
each surrogate;
\item freeze preprocessing and hyperparameters before opening an outer bout or
mouse exactly once.
\end{enumerate}

The next useful experiment is not a wider \(W,d,r\) search on these same
targets.  It is a small prespecified observable comparison using separated
NREM bouts or animals.  Reasonable candidates are a few causal rate filters
with biologically interpretable time constants and the fluorescence
observable.  If a compact planning comparison is required, the 180-s,
\(d=1,r=2\) fallback and the 45-s, \(d=1,r=4\) local-increment stress test can
be carried forward as explicitly unverified contrasts.  Neither should be
promoted.  Sliding-window Grassmann tracking should begin only after one local
estimator passes prediction, reconstruction, damping, and recurrence checks
on separated data.

\section*{Reproducibility}

The executable tutorial is
\path{DMD/scripts/02_systematic_window_tuning_mouse02_sleep.py}.  The completed
run uses PyDMD 2025.8.1 and seed 20260717.  Complete origin and horizon scores,
oracle tables, conditioning and warning tables, mode inventories and matches,
synchronized surrogate residuals, reconstruction diagnostics, figures, the
frozen decision, and a run manifest are in
\path{DMD/results/06_window_parameter_tuning/}.  The manifest records the
exact script hash used for each completed run.  The implementation contains
noise-free autonomous-forecast, independent-pair \(G,A,L\), and exact
all-neuron POD-score identity checks.  The broader DMD numerical test suite
contains 25 unit tests.

\begin{thebibliography}{9}

\bibitem{pydmdtutorial}
PyDMD developers, \emph{Tutorial 1: Dynamic Mode Decomposition on a toy
dataset}, official PyDMD 2025.8.1 documentation,
\url{https://pydmd.github.io/PyDMD/tutorial1dmd.html}.

\bibitem{hankeldmd}
PyDMD developers, \emph{Hankel DMD}, official PyDMD 2025.8.1 documentation,
\url{https://pydmd.github.io/PyDMD/hankeldmd.html}.

\bibitem{bopdmd}
PyDMD developers, \emph{BOP-DMD}, official PyDMD 2025.8.1 documentation,
\url{https://pydmd.github.io/PyDMD/bopdmd.html}.

\bibitem{windoweddmd}
H.~Zhang, C.~W.~Rowley, E.~A.~Deem, and L.~N.~Cattafesta,
"Online dynamic mode decomposition for time-varying systems,"
\emph{SIAM Journal on Applied Dynamical Systems} 18(3), 1586--1609 (2019),
\url{https://doi.org/10.1137/18M1192329}.

\bibitem{lapo}
K.~Lapo, S.~M.~Ichinaga, and J.~N.~Kutz,
"A method for unsupervised learning of coherent spatiotemporal patterns in
multiscale data," \emph{Proceedings of the National Academy of Sciences}
(2025), \url{https://doi.org/10.1073/pnas.2415786122}.

\bibitem{tashman}
L.~J.~Tashman, "Out-of-sample tests of forecasting accuracy: an analysis and
review," \emph{International Journal of Forecasting} 16(4), 437--450 (2000),
\url{https://doi.org/10.1016/S0169-2070(00)00065-0}.

\bibitem{askham}
T.~Askham and J.~N.~Kutz, "Variable projection methods for an optimized
dynamic mode decomposition," \emph{SIAM Journal on Applied Dynamical
Systems} 17(1), 380--416 (2018),
\url{https://doi.org/10.1137/M1124176}.

\bibitem{resdmd}
M.~J.~Colbrook, L.~J.~Ayton, and M.~Szőke, "Residual dynamic mode
decomposition: robust and verified Koopmanism," \emph{Journal of Fluid
Mechanics} (2023), \url{https://doi.org/10.1017/jfm.2022.1052}.

\bibitem{stochasticresdmd}
M.~J.~Colbrook, Q.~Li, R.~V.~Raut, and A.~Townsend, "Beyond expectations:
residual dynamic mode decomposition and variance for stochastic dynamical
systems," \emph{Nonlinear Dynamics} 112, 2037--2061 (2024),
\url{https://doi.org/10.1007/s11071-023-09135-w}.

\end{thebibliography}

\end{document}
